%=======================================================================
\section*{Abstract}
\addcontentsline{toc}{section}{Abstract}

This project builds a complete, reproducible pipeline that classifies brain MRI slices into
\textbf{three tumor categories --- Glioma, Meningioma, and Pituitary tumor --- using classical
(non-deep) machine learning}. Rather than learning features end-to-end with a CNN, the pipeline
extracts four complementary families of handcrafted descriptors (GLCM texture statistics, Local
Binary Patterns, Discrete Wavelet Transform sub-band statistics, and Histogram of Oriented
Gradients), augments them with first-order intensity statistics, and fuses everything into a single
feature vector.

The work is organised in two phases. \textbf{Phase~1} selects the optimal hyper-parameters of every
descriptor using three fast, classifier-free class-separability metrics (Fisher Discriminant Ratio,
Mutual Information, Davies--Bouldin Index) combined into one composite score, avoiding the cost of
training a model for every parameter combination. \textbf{Phase~2} extracts features with those
optimal settings, cleans and standardises them, and benchmarks eight classifiers (SVM, KNN, Random
Forest, XGBoost, LightGBM, Logistic Regression, MLP, ExtraTrees) with exhaustive \texttt{GridSearchCV}
tuning. A dedicated statistical-validation stage (Cohen's $\kappa$, bootstrap confidence intervals,
McNemar, Wilcoxon, Friedman + Nemenyi) then establishes whether the performance differences between
models are statistically meaningful.

\begin{notebox}[Keywords:]
brain tumor, MRI, texture analysis, GLCM, LBP, DWT, HOG, feature fusion, GridSearchCV,
statistical validation, class separability.
\end{notebox}
\newpage
